\ifdefined\gkver\else\def\gkver{student}\fi
\documentclass[\gkver]{gaokaozhenti}
\begin{document}

\chapter{2021年山东}

\begin{enumerate}
\item
%% number: 1
%% paperName: 2021年山东高考第 1 题 3 分
%% typeId: 1
%% type: 单选题
%% chapter: 原子和原子核
%% point: 天然放射性
%% method: 无
%% score: 3
%% degree: 750
%% duplicateId: 0
%% body:
在测定年代较近的湖泊沉积物形成年份时，常利用沉积物中半衰期较短的 \ce{^{210}_{82}Pb}，其衰变方程为 \ce{^{210}_{82}Pb -> ^{210}_{83}Bi + X}。以下说法正确的是 \xzanswer{A}

\fourchoices[answer=A]
{衰变方程中的 X 是电子}
{升高温度可以加快 \ce{^{210}_{82}Pb} 的衰变}
{\ce{^{210}_{82}Pb} 与 \ce{^{210}_{83}Bi} 的质量差等于衰变的质量亏损}
{方程中的 X 来自于 \ce{^{210}_{82}Pb} 内质子向中子的转化}

%% answer: A
%% memo:
\memoanswer{
	A．根据质量数守恒和电荷数守恒可知，X 是电子，A 正确；

	B．半衰期非常稳定，不受温度、压强以及该物质是单质还是化合物的影响，B 错误；

	C．\ce{^{210}_{82}Pb} 与 \ce{^{210}_{83}Bi} 和电子 X 的质量差等于衰变的质量亏损，C 错误；

	D．方程中的 X 来自于 \ce{^{210}_{82}Pb} 内中子向质子的转化，D 错误。

	故选 A。
}

\item
%% number: 2
%% paperName: 2021年山东高考第 2 题 3 分
%% typeId: 1
%% type: 单选题
%% chapter: 气体性质
%% point: 热力学第一定律
%% method: 无
%% score: 3
%% degree: 650
%% duplicateId: 0
%% body:
如图所示，密封的矿泉水瓶中，距瓶口越近水的温度越高。一开口向下、导热良好的小瓶置于矿泉水瓶中，小瓶中封闭一段空气。挤压矿泉水瓶，小瓶下沉到底部；松开后，小瓶缓慢上浮，上浮过程中，小瓶内气体 \xzanswer{B}

\onepicture[h=4.8cm]{figs/02.png}

\fourchoices[answer=B]
{内能减少}
{对外界做正功}
{增加的内能大于吸收的热量}
{增加的内能等于吸收的热量}

%% answer: B
%% memo:
\memoanswer{
	A．由于越接近矿泉水瓶口，水的温度越高，因此小瓶上浮的过程中，小瓶内温度升高，内能增加，A 错误；

	B．在小瓶上升的过程中，小瓶内气体的温度逐渐升高，压强逐渐减小，根据理想气体状态方程 $\frac{pV}{T} = C$，气体体积膨胀，对外界做正功，B 正确；

	CD．由 AB 分析，小瓶上升时，小瓶内气体内能增加，气体对外做功，根据热力学第一定律 $\Delta U = W + Q$，由于气体对外做功，因此吸收的热量大于增加的内能，CD 错误。

	故选 B。
}

\item
%% number: 3
%% paperName: 2021年山东高考第 3 题 3 分
%% typeId: 1
%% type: 单选题
%% chapter: 功和能
%% point: 动能定理
%% method: 无
%% score: 3
%% degree: 750
%% duplicateId: 0
%% body:
如图所示，粗糙程度处处相同的水平桌面上有一长为 $L$ 的轻质细杆，一端可绕竖直光滑轴 $O$ 转动，另一端与质量为 $m$ 的小木块相连。木块以水平初速度 $v_{0}$ 出发，恰好能完成一个完整的圆周运动。在运动过程中，木块所受摩擦力的大小为 \xzanswer{B}

\onepicture[w=4.0cm]{figs/03.png}

\fourchoices[answer=B]
{$\frac{mv_{0}^{2}}{2\pi L}$}
{$\frac{mv_{0}^{2}}{4\pi L}$}
{$\frac{mv_{0}^{2}}{8\pi L}$}
{$\frac{mv_{0}^{2}}{16\pi L}$}

%% answer: B
%% memo:
\memoanswer{
	在运动过程中，只有摩擦力做功，而摩擦力做功与路径有关，根据动能定理

	$-f \cdot 2\pi L = 0 - \frac{1}{2}mv_{0}^{2}$

	可得摩擦力的大小：$f = \frac{mv_{0}^{2}}{4\pi L}$

	故选 B。
}

\item
%% number: 4
%% paperName: 2021年山东高考第 4 题 3 分
%% typeId: 1
%% type: 单选题
%% chapter: 气体性质
%% point: 玻意耳定律
%% method: 无
%% score: 3
%% degree: 650
%% duplicateId: 0
%% body:
血压仪由加压气囊、臂带、压强计等构成，如图所示。加压气囊可将外界空气充入臂带，压强计示数为臂带内气体的压强高于大气压强的数值，充气前臂带内气体压强为大气压强，体积为 $V$；每次挤压气囊都能将 $60 \Ucmc$ 的外界空气充入臂带中，经 5 次充气后，臂带内气体体积变为 $5V$，压强计示数为 $150 \UmmHg$。已知大气压强等于 $750 \UmmHg$，气体温度不变。忽略细管和压强计内的气体体积。则 $V$ 等于 \xzanswer{D}

\onepicture[w=6.0cm]{figs/04.png}

\fourchoices[answer=D]
{$30 \Ucmc$}
{$40 \Ucmc$}
{$50 \Ucmc$}
{$60 \Ucmc$}

%% answer: D
%% memo:
\memoanswer{
	根据玻意耳定律可知

	$p_{0}V + 5p_{0}V_{0} = p_{1} \times 5V$

	已知 $p_{0} = 750 \UmmHg$，$V_{0} = 60 \Ucmc$，$p_{1} = 750 \UmmHg + 150 \UmmHg = 900 \UmmHg$

	代入数据整理得 $V = 60 \Ucmc$

	故选 D。
}

\item
%% number: 5
%% paperName: 2021年山东高考第 5 题 3 分
%% typeId: 1
%% type: 单选题
%% chapter: 曲线运动
%% point: 万有引力定律
%% method: 无
%% score: 3
%% degree: 750
%% duplicateId: 0
%% body:
从“玉兔”登月到“祝融”探火，我国星际探测事业实现了由地月系到行星际的跨越。已知火星质量约为月球的 9 倍，半径约为月球的 2 倍，“祝融”火星车的质量约为“玉兔”月球车的 2 倍。在着陆前，“祝融”和“玉兔”都会经历一个由着陆平台支撑的悬停过程。悬停时，“祝融”与“玉兔”所受着陆平台的作用力大小之比为 \xzanswer{B}

\onepicture[w=6.0cm]{figs/05.jpg}

\fourchoices[answer=B]
{$9:1$}
{$9:2$}
{$36:1$}
{$72:1$}

%% answer: B
%% memo:
\memoanswer{
	悬停时所受平台的作用力等于万有引力，根据 $F = G\frac{mM}{R^{2}}$

	可得 $\frac{F_\text{祝融}}{F_\text{玉兔}} = G\frac{M_\text{火}m_\text{祝融}}{R_\text{火}^{2}} : G\frac{M_\text{月}m_\text{玉兔}}{R_\text{月}^{2}} = \frac{9}{2^{2}} \times 2 = \frac{9}{2}$

	故选 B。
}

\item
%% number: 6
%% paperName: 2021年山东高考第 6 题 3 分
%% typeId: 1
%% type: 单选题
%% chapter: 电场
%% point: 带电粒子的平衡
%% method: 无
%% score: 3
%% degree: 550
%% duplicateId: 0
%% body:
如图~\ref{2021山东06a} 所示，边长为 $a$ 的正方形，四个顶点上分别固定一个电荷量为 $+q$ 的点电荷；在 $0 \leqslant x < \frac{\sqrt{2}}{2}a$ 区间，$x$ 轴上电势 $\varphi$ 的变化曲线如图~\ref{2021山东06b} 所示。现将一电荷量为 $-Q$ 的点电荷 P 置于正方形的中心 $O$ 点，此时每个点电荷所受库仑力的合力均为零。若将 P 沿 $x$ 轴向右略微移动后，由静止释放，以下判断正确的是 \xzanswer{C}

\twopicture[labelA=2021山东06a, labelB=2021山东06b, numA=甲, numB=乙, wA=4.5cm, wB=3.8cm]
{figs/06a.png}
{figs/06b.png}

\fourchoices[answer=C]
{$Q = \frac{\sqrt{2}+1}{2}q$，释放后 P 将向右运动}
{$Q = \frac{\sqrt{2}+1}{2}q$，释放后 P 将向左运动}
{$Q = \frac{2\sqrt{2}+1}{4}q$，释放后 P 将向右运动}
{$Q = \frac{2\sqrt{2}+1}{4}q$，释放后 P 将向左运动}

%% answer: C
%% memo:
\memoanswer{
	对 $y$ 轴正向的点电荷，由平衡知识可得

	$\sqrt{2}k\frac{q^{2}}{a^{2}} + k\frac{q^{2}}{(\sqrt{2}a)^{2}} = k\frac{Qq}{(\frac{\sqrt{2}}{2}a)^{2}}$

	解得：$Q = \frac{2\sqrt{2}+1}{4}q$

	因在 $0 \leqslant x < \frac{\sqrt{2}}{2}a$ 区间内沿 $x$ 轴正向电势升高，则场强方向沿 $x$ 轴负向，则将 P 沿 $x$ 轴正向向右略微移动后释放，P 受到向右的电场力而向右运动。

	故选 C。
}

\item
%% number: 7
%% paperName: 2021年山东高考第 7 题 3 分
%% typeId: 1
%% type: 单选题
%% chapter: 光学
%% point: 光的干涉
%% method: 无
%% score: 3
%% degree: 450
%% duplicateId: 0
%% body:
用平行单色光垂直照射一层透明薄膜，观察到如图所示明暗相间的干涉条纹。下列关于该区域薄膜厚度 $d$ 随坐标 $x$ 的变化图像，可能正确的是 \xzanswer{D}

\onepicture[w=7.0cm]{figs/07a.png}

\fourchoices[answer=D, ispicture=true, w=2.8cm]
{figs/07b.png}
{figs/07c.png}
{figs/07d.png}
{figs/07e.png}

%% answer: D
%% memo:
\memoanswer{
	从薄膜的上下表面分别反射的两列光是相干光，其光程差为 $\Delta x = 2d$，即光程差为薄膜厚度的 2 倍，当光程差 $\Delta x = n\lambda$ 时此处表现为亮条纹，故相邻亮条纹之间的薄膜的厚度差为 $\frac{1}{2}\lambda$，在图中相邻亮条纹（或暗条纹）之间的距离变大，则薄膜层的厚度之间变小，因条纹宽度逐渐变宽，则厚度不是均匀变小。

	故选 D。
}

\item
%% number: 8
%% paperName: 2021年山东高考第 8 题 3 分
%% typeId: 1
%% type: 单选题
%% chapter: 电磁感应
%% point: 切割磁感线电动势
%% method: 无
%% score: 3
%% degree: 450
%% duplicateId: 0
%% body:
迷你系绳卫星在地球赤道正上方的电离层中，沿圆形轨道绕地飞行。系绳卫星由两子卫星组成，它们之间的导体绳沿地球半径方向，如图所示。在电池和感应电动势的共同作用下，导体绳中形成指向地心的电流，等效总电阻为 $r$。导体绳所受的安培力克服大小为 $f$ 的环境阻力，可使卫星保持在原轨道上。已知卫星离地平均高度为 $H$，导体绳长为 $L$（$L \ll H$），地球半径为 $R$，质量为 $M$，轨道处磁感应强度大小为 $B$，方向垂直于赤道平面。忽略地球自转的影响。据此可得，电池电动势为 \xzanswer{A}

\onepicture[h=5.5cm]{figs/08.png}

\fourchoices[answer=A]
{$BL\sqrt{\frac{GM}{R+H}} + \frac{fr}{BL}$}
{$BL\sqrt{\frac{GM}{R+H}} - \frac{fr}{BL}$}
{$BL\sqrt{\frac{GM}{R+H}} + \frac{BL}{fr}$}
{$BL\sqrt{\frac{GM}{R+H}} - \frac{BL}{fr}$}

%% answer: A
%% memo:
\memoanswer{
	根据 $G\frac{Mm}{(R+H)^{2}} = m\frac{v^{2}}{R+H}$

	可得卫星做圆周运动的线速度 $v = \sqrt{\frac{GM}{R+H}}$

	根据右手定则可知，导体绳产生的感应电动势相当于上端为正极的电源，其大小为 $E' = BLv$

	因导线绳所受阻力 $f$ 与安培力 $F$ 平衡，则安培力与速度方向相同，可知导线绳中的电流方向向下，即电池电动势大于导线绳切割磁感线产生的电动势，可得 $f = B\frac{E - E'}{r}L$

	解得 $E = BL\sqrt{\frac{GM}{R+H}} + \frac{fr}{BL}$

	故选 A。
}

\item
%% number: 9
%% paperName: 2021年山东高考第 9 题 4 分
%% typeId: 2
%% type: 多选题
%% chapter: 交流电
%% point: 变压器
%% method: 无
%% score: 4
%% degree: 650
%% duplicateId: 0
%% body:
输电能耗演示电路如图所示。左侧变压器原、副线圈匝数比为 $1:3$，输入电压为 $7.5 \UV$ 的正弦交流电。连接两理想变压器的导线总电阻为 $r$，负载 $R$ 的阻值为 $10 \UO$。开关 S 接 1 时，右侧变压器原、副线圈匝数比为 $2:1$，$R$ 上的功率为 $10 \UW$；接 2 时，匝数比为 $1:2$，$R$ 上的功率为 $P$。以下判断正确的是 \xzanswer{BD}

\onepicture[w=9.0cm]{figs/09.png}

\fourchoices[answer=BD]
{$r = 10 \UO$}
{$r = 5 \UO$}
{$P = 45 \UW$}
{$P = 22.5 \UW$}

%% answer: BD
%% memo:
\memoanswer{
	当开关 S 接 1 时，左侧变压器次级电压 $U_{2} = 3 \times 7.5 \UV = 22.5 \UV$

	电阻 $R$ 上的电压，即右侧变压器的次级电压 $U_{4} = \sqrt{PR} = 10 \UV$

	电流 $I_{4} = \frac{U_{4}}{R} = 1 \UA$

	则右侧变压器初级电压 $U_{3} = \frac{2}{1} \times 10 \UV = 20 \UV$

	电流 $I_{3} = \frac{1}{2} \times 1 \UA = 0.5 \UA$

	则 $r = \frac{U_{2} - U_{3}}{I_{3}} = 5 \UO$

	当开关 S 接 2 时，设输电电流为 $I$，则右侧变压器的次级电流为 $0.5I$；右侧变压器两边电压关系可知

	$\frac{U_{2} - Ir}{n_{3}} = \frac{0.5IR}{n_{4}}$

	解得 $I = 3 \UA$

	则 $R$ 上的功率 $P = (0.5I)^{2}R = 22.5 \UW$

	故选 BD。
}

\item
%% number: 10
%% paperName: 2021年山东高考第 10 题 4 分
%% typeId: 2
%% type: 多选题
%% chapter: 机械波
%% point: 波与振动图象
%% method: 无
%% score: 4
%% degree: 550
%% duplicateId: 0
%% body:
一列简谐横波沿 $x$ 轴传播，如图所示，实线为 $t_{1} = 2 \Us$ 时的波形图，虚线为 $t_{2} = 5 \Us$ 时的波形图。以下关于平衡位置在 $O$ 处质点的振动图像，可能正确的是 \xzanswer{AC}

\onepicture[w=7.0cm]{figs/10a.png}

\fourchoices[answer=AC, ispicture=true, w=3.4cm]
{figs/10b.png}
{figs/10c.png}
{figs/10d.png}
{figs/10e.png}

%% answer: AC
%% memo:
\memoanswer{
	机械波的传播方向不确定，所以需要考虑机械波传播方向的不确定性。

	AB．若机械波沿 $x$ 轴正方向传播，在 $t_{1} = 2 \Us$ 时 $O$ 点振动方向竖直向上，则传播时间 $\Delta t = t_{2} - t_{1} = 3 \Us$ 满足

	$\Delta t = \frac{3}{4}T + nT$（$n = 0,1,2,3\cdots$）

	解得 $T = \frac{12}{4n+3} \Us$（$n = 0,1,2,3\cdots$）

	当 $n = 0$ 时，解得周期 $T = 4 \Us$。

	$A$ 正确，$B$ 错误；

	CD．若机械波沿 $x$ 轴负方向传播，在 $t_{2} = 5 \Us$ 时 $O$ 点处于波谷，则

	$\Delta t = \frac{1}{4}T + nT$（$n = 0,1,2,3\cdots$）

	解得 $T = \frac{12}{4n+1} \Us$（$n = 0,1,2,3\cdots$）

	当 $n = 0$ 时，解得周期 $T = 12 \Us$。

	$C$ 正确，$D$ 错误。

	故选 AC。
}

\item
%% number: 11
%% paperName: 2021年山东高考第 11 题 4 分
%% typeId: 2
%% type: 多选题
%% chapter: 曲线运动
%% point: 平抛运动
%% method: 无
%% score: 4
%% degree: 450
%% duplicateId: 0
%% body:
如图所示，载有物资的热气球静止于距水平地面 $H$ 的高处，现将质量为 $m$ 的物资以相对地面的速度 $v_{0}$ 水平投出，落地时物资与热气球的距离为 $d$。已知投出物资后热气球的总质量为 $M$，所受浮力不变，重力加速度为 $g$，不计阻力，以下判断正确的是 \xzanswer{BC}

\onepicture[h=5.0cm]{figs/11.png}

\fourchoices[answer=BC]
{投出物资后热气球做匀加速直线运动}
{投出物资后热气球所受合力大小为 $mg$}
{$d = \left(1 + \frac{m}{M}\right)\sqrt{\frac{2Hv_{0}^{2}}{g} + H^{2}}$}
{$d = \sqrt{\frac{2Hv_{0}^{2}}{g} + \left(1 + \frac{m}{M}\right)^{2}H^{2}}$}

%% answer: BC
%% memo:
\memoanswer{
	AB．热气球开始携带物资时处于静止状态，所受合外力为 0，初动量为 0，水平投出重力为 $mg$ 的物资瞬间，满足动量守恒定律 $Mv = mv_{0}$

	则热气球和物资的动量等大反向，热气球获得水平向左的速度 $v$，热气球所受合外力恒为 $mg$，竖直向上，所以热气球做匀加速曲线运动，A 错误，B 正确；

	CD．热气球和物资的运动示意图如图所示

	\onepicture[w=6.0cm]{figs/11_an.jpg}

	热气球和物资所受合力大小均为 $mg$，所以热气球在竖直方向上加速度大小为 $a = \frac{m}{M}g$

	物资落地 $H$ 过程所用的时间 $t$ 内，根据 $H = \frac{1}{2}gt^{2}$ 解得落地时间为 $t = \sqrt{\frac{2H}{g}}$

	热气球在竖直方向上运动的位移为 $H_{M} = \frac{1}{2}at^{2} = \frac{1}{2} \cdot \frac{m}{M}g \cdot \frac{2H}{g} = \frac{m}{M}H$

	热气球和物资在水平方向均做匀速直线运动，水平位移为 $x_{m} = v_{0}t = v_{0}\sqrt{\frac{2H}{g}}$，$x_{M} = vt = \frac{m}{M}v_{0} \cdot \sqrt{\frac{2H}{g}}$

	根据勾股定理可知热气球和物资的实际位移为 $d = \sqrt{(x_{m} + x_{M})^{2} + (H + H_{M})^{2}} = \left(1 + \frac{m}{M}\right)\sqrt{\frac{2Hv_{0}^{2}}{g} + H^{2}}$

	C 正确，D 错误。

	故选 BC。
}

\item
%% number: 12
%% paperName: 2021年山东高考第 12 题 4 分
%% typeId: 2
%% type: 多选题
%% chapter: 电磁感应
%% point: 电磁感应中的变加速
%% method: 无
%% score: 4
%% degree: 350
%% duplicateId: 0
%% body:
如图所示，电阻不计的光滑 U 形金属导轨固定在绝缘斜面上。区域Ⅰ、Ⅱ中磁场方向均垂直斜面向上，Ⅰ区中磁感应强度随时间均匀增加，Ⅱ区中为匀强磁场。阻值恒定的金属棒从无磁场区域中 a 处由静止释放，进入Ⅱ区后，经 b 下行至 c 处反向上行。运动过程中金属棒始终垂直导轨且接触良好。在第一次下行和上行的过程中，以下叙述正确的是 \xzanswer{ABD}

\onepicture[w=6.0cm]{figs/12.png}

\fourchoices[answer=ABD]
{金属棒下行过 b 时的速度大于上行过 b 时的速度}
{金属棒下行过 b 时的加速度大于上行过 b 时的加速度}
{金属棒不能回到无磁场区}
{金属棒能回到无磁场区，但不能回到 a 处}

%% answer: ABD
%% memo:
\memoanswer{
	AB．在Ⅰ区域中，磁感应强度为 $B_{1} = kt$，感应电动势 $E_{1} = \frac{\Delta B_{1}}{\Delta t}S = kS$

	感应电动势恒定，所以导体棒上的感应电流恒为 $I_{1} = \frac{E_{1}}{R} = \frac{kS}{R}$

	导体棒进入Ⅱ区域后，导体切割磁感线，产生一个感应电动势，因为导体棒到达 c 点后又能上行，说明加速度始终沿斜面向上，下行和上行经过 b 点的受力分析如图

	\onepicture[w=13.0cm]{figs/12_an.jpg}

	设下行、上行过 b 时导体棒的速度分别为 $v$、$v'$，则下行过 b 时导体棒切割磁感线产生的感应电流为 $E_{2} = B_{2}Lv$

	下行过 b 时导体棒上的电流为 $I_{2} = \frac{E_{2} + E_{1}}{R} = \frac{B_{2}Lv}{R} + \frac{kS}{R}$

	下行过 b 时，根据牛顿第二定律可知 $B_{2}I_{2}L - mg\sin\theta = \frac{B_{2}^{2}L^{2}v}{R} + \frac{B_{2}kSL}{R} - mg\sin\theta = ma_{1}$

	上行过 b 时，切割磁感线产生的感应电动势为 $E_{2}' = B_{2}Lv'$

	上行过 b 时导体棒上的电流为 $I_{3} = \frac{E_{1} - E_{2}'}{R} = \frac{kS}{R} - \frac{B_{2}Lv'}{R}$

	根据牛顿第二定律可知 $B_{2}I_{3}L - mg\sin\theta = \frac{B_{2}kSL}{R} - \frac{B_{2}^{2}L^{2}v'}{R} - mg\sin\theta = ma_{2}$

	比较加速度大小可知 $a_{1} > a_{2}$

	由于 bc 段距离不变，下行过程中加速度大，上行过程中加速度小，所以金属棒下行经过 b 点时的速度大于上行经过 b 点时的速度，AB 正确；

	CD．导体棒上行时，加速度与速度同向，则导体棒做加速度减小的加速运动，则一定能回到无磁场区。由 AB 分析可得，导体棒进磁场Ⅱ区（下行进磁场）的速度大于出磁场Ⅱ区（下行进磁场）的速度，导体棒在无磁场区做加速度相同的减速运动 $mg\sin\theta = ma_{0}$

	则金属棒不能回到 a 处，C 错误，D 正确。

	故选 ABD。
}

\item
%% number: 13
%% paperName: 2021年山东高考第 13 题 6 分
%% typeId: 4
%% type: 实验题
%% chapter: 直线运动
%% point: 竖直上下抛运动
%% method: 无
%% score: 6
%% degree: 550
%% duplicateId: 0
%% body:
某乒乓球爱好者，利用手机研究乒乓球与球台碰撞过程中能量损失的情况。实验步骤如下：

\begin{steps}
	\item
	固定好手机，打开录音功能；
	\item
	从一定高度由静止释放乒乓球；
	\item
	手机记录下乒乓球与台面碰撞的声音，其随时间（单位：$\Us[a]$）的变化图像如图所示。
\end{steps}

根据声音图像记录的碰撞次序及相应碰撞时刻，如下表所示。

\begin{table}[htbp]
\centering
\begin{tblr}{|c|c|c|c|c|c|c|c|}
\hline
碰撞次序 & 1 & 2 & 3 & 4 & 5 & 6 & 7 \\
\hline
碰撞时刻（$\Us[a]$） & 1.12 & 1.58 & 2.00 & 2.40 & 2.78 & 3.14 & 347 \\
\hline
\end{tblr}
\end{table}

根据实验数据，回答下列问题：
\begin{enumerate}
	\item
	利用碰撞时间间隔，计算出第 3 次碰撞后乒乓球的弹起高度为 \tkanswer{0.20} $\Um$（保留 2 位有效数字，当地重力加速度 $g = 9.80 \Umsq$）。
	\item
	设碰撞后弹起瞬间与该次碰撞前瞬间速度大小的比值为 $k$，则每次碰撞损失的动能为碰撞前动能的 \tkanswer{$1-k^{2}$} 倍（用 $k$ 表示），第 3 次碰撞过程中 $k =$ \tkanswer{0.95}（保留 2 位有效数字）。
	\item
	由于存在空气阻力，第（1）问中计算的弹起高度 \tkanswer{高于}（填“高于”或“低于”）实际弹起高度。
\end{enumerate}

\onepicture[w=12.0cm, align=right]{figs/13.png}

%% answer:
\jdanswer{
\begin{enumerate}
	\item
	$0.20$
	\item
	$1-k^{2}$，$0.95$
	\item
	高于
\end{enumerate}
}
%% memo:
\memoanswer{
\begin{enumerate}
	\item
	第 3 次碰撞到第 4 次碰撞用时 $t_{0} = (2.40 - 2.00) \Us = 0.40 \Us$，根据竖直上抛和自由落体运动的对称性可知第 3 次碰撞后乒乓球弹起的高度为

	$h_{0} = \frac{1}{2}g\left(\frac{t_{0}}{2}\right)^{2} = \frac{1}{2} \times 9.8 \times 0.2^{2} \Um \approx 0.20 \Um$

	\item
	碰撞后弹起瞬间速度为 $v_{2}$，碰撞前瞬间速度为 $v_{1}$，根据题意可知 $\frac{v_{2}}{v_{1}} = k$

	则每次碰撞损失的动能与碰撞前动能的比值为 $\frac{\frac{1}{2}mv_{1}^{2} - \frac{1}{2}mv_{2}^{2}}{\frac{1}{2}mv_{1}^{2}} = 1 - \frac{\frac{1}{2}mv_{2}^{2}}{\frac{1}{2}mv_{1}^{2}} = 1 - k^{2}$

	第 2 次碰后从最高点落地瞬间的速度 $v = gt = \left(\frac{2.00 - 1.58}{2}\right)g = 0.21g$

	第 3 次碰撞后瞬间速度为 $v' = gt' = \left(\frac{2.40 - 2.00}{2}\right)g = 0.20g$

	则第 3 次碰撞过程中 $k = \frac{v'}{v} = \frac{0.20}{0.21} \approx 0.95$

	\item
	由于存在空气阻力，乒乓球在上升过程中受到向下的阻力和重力，加速度变大，上升的高度变小，所以第（1）问中计算的弹起高度高于实际弹起的高度。
\end{enumerate}
}

\item
%% number: 14
%% paperName: 2021年山东高考第 14 题 8 分
%% typeId: 4
%% type: 实验题
%% chapter: 电路
%% point: 传感器
%% method: 无
%% score: 8
%% degree: 450
%% duplicateId: 0
%% body:
热敏电阻是传感器中经常使用的元件，某学习小组要探究一热敏电阻的阻值随温度变化的规律。可供选择的器材有：
\begin{itemize}
	\item
	待测热敏电阻 $R_{T}$（实验温度范围内，阻值约几百欧到几千欧）；
	\item
	电源 $E$（电动势 $1.5 \UV$，内阻 $r$ 约为 $0.5 \UO$）；
	\item
	电阻箱 $R$（阻值范围 $0 \sim 9999.99 \UO$）；
	\item
	滑动变阻器 $R_{1}$（最大阻值 $20 \UO$）；
	\item
	滑动变阻器 $R_{2}$（最大阻值 $2000 \UO$）；
	\item
	微安表（量程 $100 \UuA$，内阻等于 $2500 \UO$）；
	\item
	开关两个，温控装置一套，导线若干。
\end{itemize}

同学们设计了如图~\ref{2021山东14a} 所示的测量电路，主要实验步骤如下：
\begin{steps}
	\item
	按图示连接电路；
	\item
	闭合 $S_{1}$、$S_{2}$，调节滑动变阻器滑片 P 的位置，使微安表指针满偏；
	\item
	保持滑动变阻器滑片 P 的位置不变，断开 $S_{2}$，调节电阻箱，使微安表指针半偏；
	\item
	记录此时的温度和电阻箱的阻值。
\end{steps}

回答下列问题：
\begin{enumerate}
	\item
	为了更准确地测量热敏电阻的阻值，滑动变阻器应选用 \tkanswer{$R_{1}$}（填“$R_{1}$”或“$R_{2}$”）。
	\item
	请用笔画线代替导线，将实物图（不含温控装置）连接成完整电路 \tkanswer[3cm]{}。
	\item
	某温度下微安表半偏时，电阻箱的读数为 $6000.00 \UO$，该温度下热敏电阻的测量值为 \tkanswer{3500} $\UO$（结果保留到个位），该测量值 \tkanswer{大于}（填“大于”或“小于”）真实值。
	\item
	多次实验后，学习小组绘制了如图~\ref{2021山东14c} 所示的图像。由图像可知，该热敏电阻的阻值随温度的升高逐渐 \tkanswer{减小}（填“增大”或“减小”）。
\end{enumerate}

\onepicture[num=甲, label=2021山东14a, align=right, w=6.5cm]{figs/14a.png}
\onepicture[align=right, w=5.0cm]{figs/14b.png}
\onepicture[num=乙, label=2021山东14c, align=right, w=6.5cm]{figs/14c.png}

%% answer:
\jdanswer{
\begin{enumerate}
	\item
	$R_{1}$
	\item
	连线如图所示
	\item
	$3500$，大于
	\item
	减小
\end{enumerate}
}
%% memo:
\memoanswer{
（1）用半偏法测量热敏电阻的阻值，尽可能让该电路的电压在 $S_{2}$ 闭合前、后保持不变，由于该支路与滑动变阻器左侧部分电阻并联，滑动变阻器的阻值越小，$S_{2}$ 闭合前、后并联部分电阻变化越小，从而并联部分的电压值变化越小，故滑动变阻器应选 $R_{1}$。

（2）电路连接图如图所示

\onepicture[w=6.0cm]{figs/14_an.png}

（3）微安表半偏时，该支路的总电阻为原来的 2 倍，即 $R_{T} + R_{\mu A} = R = 6000.00 \UO$

可得 $R_{T} = 3500 \UO$

当断开 $S_{2}$，微安表半偏时，由于该支路的电阻增加，电压略有升高，根据欧姆定律，总电阻比原来 2 倍略大，也就是电阻箱的阻值略大于热敏电阻与微安表的总电阻，而我们用电阻箱的阻值等于热敏电阻与微安表的总电阻来计算，因此热敏电阻的测量值比真实值偏大。

（4）由于是 $\ln R_{T} - \frac{1}{T}$ 图像，当温度 $T$ 升高时，$\frac{1}{T}$ 减小，从图中可以看出 $\ln R_{T}$ 减小，从而 $R_{T}$ 减小，因此热敏电阻随温度的升高逐渐减小。
}

\item
%% number: 15
%% paperName: 2021年山东高考第 15 题 7 分
%% typeId: 6
%% type: 计算题
%% chapter: 光学
%% point: 光的折射
%% method: 无
%% score: 7
%% degree: 550
%% duplicateId: 0
%% body:
超强超短光脉冲产生方法曾获诺贝尔物理学奖，其中用到的一种脉冲激光展宽器截面如图所示。在空气中对称放置四个相同的直角三棱镜，顶角为 $\theta$。一细束脉冲激光垂直第一个棱镜左侧面入射，经过前两个棱镜后分为平行的光束，再经过后两个棱镜重新合成为一束，此时不同频率的光前后分开，完成脉冲展宽。已知相邻两棱镜斜面间的距离 $d = 100.0 \Umm$，脉冲激光中包含两种频率的光，它们在棱镜中的折射率分别为 $n_{1} = \sqrt{2}$ 和 $n_{2} = \frac{\sqrt{31}}{4}$。取 $\sin37 = 0.6$，$\cos37^\circ = 0.8$，$\frac{5}{\sqrt{7}} = 1.890$。
\begin{enumerate}
	\item
	为使两种频率的光都能从左侧第一个棱镜斜面射出，求 $\theta$ 的取值范围；
	\item
	若 $\theta = 37^\circ$，求两种频率的光通过整个展宽器的过程中，在空气中的路程差 $\Delta L$（保留 3 位有效数字）。
\end{enumerate}

\onepicture[align=right, w=7.0cm]{figs/15.png}

%% answer:
\jdanswer{
\begin{enumerate}
	\item
	$0 < \theta < 45^\circ$（或 $\theta < 45^\circ$）
	\item
	$\Delta L = 14.4 \Umm$
\end{enumerate}
}
%% memo:
\memoanswer{
（1）由几何关系可得，光线在第一个三棱镜右侧斜面上的入射角等于 $\theta$，要使得两种频率的光都从左侧第一个棱镜斜面射出，则 $\theta$ 需要比两种频率光线的全反射角都小，设 $C$ 是全反射的临界角，根据折射定律得

$\sin C = \frac{1}{n}$ ①

折射率越大，临界角越小，代入较大的折射率得

$C_{\min} = 45^\circ$ ②

所以顶角 $\theta$ 的范围为

$0 < \theta < 45^\circ$（或 $\theta < 45^\circ$）③

（2）脉冲激光从第一个三棱镜右侧斜面射出时发生折射，设折射角分别为 $\alpha_{1}$ 和 $\alpha_{2}$，由折射定律得

$n_{1} = \frac{\sin\alpha_{1}}{\sin\theta}$ ④

$n_{2} = \frac{\sin\alpha_{2}}{\sin\theta}$ ⑤

设两束光在前两个三棱镜斜面之间的路程分别为 $L_{1}$ 和 $L_{2}$，则

$L_{1} = \frac{d}{\cos\alpha_{1}}$ ⑥

$L_{2} = \frac{d}{\cos\alpha_{2}}$ ⑦

$\Delta L = 2(L_{1} - L_{2})$ ⑧

联立④⑤⑥⑦⑧式，代入数据得

$\Delta L = 14.4 \Umm$ ⑨
}

\item
%% number: 16
%% paperName: 2021年山东高考第 16 题 9 分
%% typeId: 6
%% type: 计算题
%% chapter: 曲线运动
%% point: 平抛运动
%% method: 无
%% score: 9
%% degree: 550
%% duplicateId: 0
%% body:
海鸥捕到外壳坚硬的鸟蛤（贝类动物）后，有时会飞到空中将它丢下，利用地面的冲击打碎硬壳。一只海鸥叼着质量 $m = 0.1 \Ug$ 的鸟蛤，在 $H = 20 \Um$ 的高度、以 $v_{0} = 15 \Ums$ 的水平速度飞行时，松开嘴巴让鸟蛤落到水平地面上。取重力加速度 $g = 10 \Umsq$，忽略空气阻力。
\begin{enumerate}
	\item
	若鸟蛤与地面的碰撞时间 $\Delta t = 0.005 \Us$，弹起速度可忽略，求碰撞过程中鸟蛤受到的平均作用力的大小 $F$；（碰撞过程中不计重力）
	\item
	在海鸥飞行方向正下方的地面上，有一与地面平齐、长度 $L = 6 \Um$ 的岩石，以岩石左端为坐标原点，建立如图所示坐标系。若海鸥水平飞行的高度仍为 $20 \Um$，速度大小在 $15 \Ums \sim 17 \Ums$ 之间，为保证鸟蛤一定能落到岩石上，求释放鸟蛤位置的 $x$ 坐标范围。
\end{enumerate}

\onepicture[align=right, w=6.0cm]{figs/16.png}

%% answer:
\jdanswer{
\begin{enumerate}
	\item
	$F = 500 \UN$
	\item
	$[34 \Um, 36 \Um]$ 或 $(34 \Um, 36 \Um)$
\end{enumerate}
}
%% memo:
\memoanswer{
（1）设平抛运动的时间为 $t$，鸟蛤落地前瞬间的速度大小为 $v$。竖直方向分速度大小为 $v_{y}$，根据运动的合成与分解得

$H = \frac{1}{2}gt^{2}, \quad v_{y} = gt, \quad v = \sqrt{v_{0}^{2} + v_{y}^{2}}$

在碰撞过程中，以鸟蛤为研究对象，取速度 $v$ 的方向为正方向，由动量定理得

$-F\Delta t = 0 - mv$

联立，代入数据得 $F = 500 \UN$

（2）若释放鸟蛤的初速度为 $v_{1} = 15 \Ums$，设击中岩石左端时，释放点的 $x$ 坐标为 $x_{1}$，击中右端时，释放点的 $x$ 坐标为 $x_{2}$，得

$x_{1} = v_{1}t, \quad x_{2} = x_{1} + L$

联立，代入数据得 $x_{1} = 30 \Um, \quad x_{2} = 36 \Um$

若释放鸟蛤时的初速度为 $v_{2} = 17 \Ums$，设击中岩石左端时，释放点的 $x$ 坐标为 $x_{1}'$，击中右端时，释放点的 $x$ 坐标为 $x_{2}'$，得

$x_{1}' = v_{2}t, \quad x_{2}' = x_{1}' + L$

联立，代入数据得 $x_{1}' = 34 \Um, \quad x_{2}' = 40 \Um$

综上得 $x$ 坐标区间

$[34 \Um, 36 \Um]$ 或 $(34 \Um, 36 \Um)$
}

\item
%% number: 17
%% paperName: 2021年山东高考第 17 题 14 分
%% typeId: 6
%% type: 计算题
%% chapter: 磁场
%% point: 带电粒子在磁场中的运动
%% method: 无
%% score: 14
%% degree: 350
%% duplicateId: 0
%% body:
某离子实验装置的基本原理如图~\ref{2021山东17a} 所示。Ⅰ区宽度为 $d$，左边界与 $x$ 轴垂直交于坐标原点 $O$，其内充满垂直于 $xOy$ 平面向里的匀强磁场，磁感应强度大小为 $B_{0}$；Ⅱ区宽度为 $L$，左边界与 $x$ 轴垂直交于 $O_{1}$ 点，右边界与 $x$ 轴垂直交于 $O_{2}$ 点，其内充满沿 $y$ 轴负方向的匀强电场。测试板垂直 $x$ 轴置于Ⅱ区右边界，其中心 $C$ 与 $O_{2}$ 点重合。从离子源不断飘出电荷量为 $q$、质量为 $m$ 的正离子，加速后沿 $x$ 轴正方向过 $O$ 点，依次经Ⅰ区、Ⅱ区，恰好到达测试板中心 $C$。已知离子刚进入Ⅱ区时速度方向与 $x$ 轴正方向的夹角为 $\theta$。忽略离子间的相互作用，不计重力。
\begin{enumerate}
	\item
	求离子在Ⅰ区中运动时速度的大小 $v$；
	\item
	求Ⅱ区内电场强度的大小 $E$；
	\item
	保持上述条件不变，将Ⅱ区分为左右两部分，分别填充磁感应强度大小均为 $B$（数值未知）方向相反且平行 $y$ 轴的匀强磁场，如图~\ref{2021山东17b} 所示。为使离子的运动轨迹与测试板相切于 $C$ 点，需沿 $x$ 轴移动测试板，求移动后 $C$ 到 $O_{1}$ 的距离 $S$。
\end{enumerate}

\twopicture[labelA=2021山东17a, labelB=2021山东17b, numA=甲, numB=乙, wA=6.0cm, wB=6.0cm, align=right]
{figs/17a.png}
{figs/17b.png}

%% answer:
\jdanswer{
\begin{enumerate}
	\item
	$v = \frac{qB_{0}d}{m\sin\theta}$
	\item
	$E = \frac{2qB_{0}^{2}d^{2}}{mL^{2}\tan^{2}\theta}\left(L\tan\theta + \frac{d}{\sin\theta} - \frac{d}{\tan\theta}\right)$
	\item
	$S = \frac{6(\sqrt{3} + 1)}{7\pi}L$
\end{enumerate}
}
%% memo:
\memoanswer{
（1）设离子在Ⅰ区内做匀速圆周运动的半径为 $r$，由牛顿第二定律得

$qvB_{0} = m\frac{v^{2}}{r}$ ①

根据几何关系得

$\sin\theta = \frac{d}{r}$ ②

联立①②式得

$v = \frac{qB_{0}d}{m\sin\theta}$

（2）离子在Ⅱ区内只受电场力，$x$ 方向做匀速直线运动，$y$ 方向做匀变速直线运动，设从进入电场到击中测试板中心 $C$ 的时间为 $t$，$y$ 方向的位移为 $y_{0}$，加速度大小为 $a$，由牛顿第二定律得

$qE = ma$

由运动的合成与分解得

$L = (v\cos\theta)t, \quad y_{0} = -r(1 - \cos\theta), \quad y_{0} = (v\sin\theta)t - \frac{1}{2}at^{2}$

联立得

$E = \frac{2qB_{0}^{2}d^{2}}{mL^{2}\tan^{2}\theta}\left(L\tan\theta + \frac{d}{\sin\theta} - \frac{d}{\tan\theta}\right)$

（3）Ⅱ区内填充磁场后，离子在垂直 $y$ 轴的方向做线速度大小为 $v\cos\theta$ 的匀速圆周运动，如图所示。设左侧部分的圆心角为 $\alpha$，圆周运动半径为 $r'$，运动轨迹长度为 $l'$，由几何关系得

$l' = \frac{\alpha}{2\pi} \times 2\pi r' + \frac{\alpha + \frac{\pi}{2}}{2\pi} \times 2\pi r', \quad \cos\alpha = \frac{r'}{2r'}$

由于在 $y$ 轴方向的运动不变，离子的运动轨迹与测试板相切于 $C$ 点，则离子在Ⅱ区内的运动时间不变，故有

$\frac{l'}{v\cos\theta} = \frac{L}{v\cos\theta}$

$C$ 到 $O_{1}$ 的距离

$S = 2r'\sin\alpha + r'$

联立得

$S = \frac{6(\sqrt{3} + 1)}{7\pi}L$

\onepicture[w=6.0cm]{figs/17_an.jpg}
}

\item
%% number: 18
%% paperName: 2021年山东高考第 18 题 16 分
%% typeId: 6
%% type: 计算题
%% chapter: 功和能
%% point: 动能定理
%% method: 无
%% score: 16
%% degree: 250
%% duplicateId: 0
%% body:
如图所示，三个质量均为 $m$ 的小物块 A、B、C，放置在水平地面上，A 紧靠竖直墙壁，一劲度系数为 $k$ 的轻弹簧将 A、B 连接，C 紧靠 B，开始时弹簧处于原长，A、B、C 均静止。现给 C 施加一水平向左、大小为 $F$ 的恒力，使 B、C 一起向左运动，当速度为零时，立即撤去恒力，一段时间后 A 离开墙壁，最终三物块都停止运动。已知 A、B、C 与地面间的滑动摩擦力大小均为 $f$，最大静摩擦力等于滑动摩擦力，弹簧始终在弹性限度内。（弹簧的弹性势能可表示为：$E_{p} = \frac{1}{2}kx^{2}$，$k$ 为弹簧的劲度系数，$x$ 为弹簧的形变量）
\begin{enumerate}
	\item
	求 B、C 向左移动的最大距离 $x_{0}$ 和 B、C 分离时 B 的动能 $E_{k}$；
	\item
	为保证 A 能离开墙壁，求恒力的最小值 $F_{\min}$；
	\item
	若三物块都停止时 B、C 间的距离为 $x_{BC}$，从 B、C 分离到 B 停止运动的整个过程，B 克服弹簧弹力做的功为 $W$，通过推导比较 $W$ 与 $fx_{BC}$ 的大小；
	\item
	若 $F = 5f$，请在所给坐标系中，画出 C 向右运动过程中加速度 $a$ 随位移 $x$ 变化的图像，并在坐标轴上标出开始运动和停止运动时的 $a$、$x$ 值（用 $f$、$k$、$m$ 表示），不要求推导过程。以撤去 $F$ 时 C 的位置为坐标原点，水平向右为正方向。
\end{enumerate}

\onepicture[align=right, w=10.0cm]{figs/18.png}

%% answer:
\jdanswer{
\begin{enumerate}
	\item
	$x_{0} = \frac{2F - 4f}{k}$，$E_{k} = \frac{F^{2} - 6fF + 8f^{2}}{k}$
	\item
	$F_{\min} = \left(3 + \frac{\sqrt{10}}{2}\right)f$
	\item
	$W < fx_{BC}$
	\item
	图象如图所示
\end{enumerate}
}
%% memo:
\memoanswer{
（1）从开始到 B、C 向左移动到最大距离的过程中，以 B、C 和弹簧为研究对象，由功能关系得

$Fx_{0} = 2fx_{0} + \frac{1}{2}kx_{0}^{2}$

弹簧恢复原长时 B、C 分离，从弹簧最短到 B、C 分离，以 B、C 和弹簧为研究对象，由能量守恒得

$\frac{1}{2}kx_{0}^{2} = 2fx_{0} + 2E_{k}$

联立方程解得

$x_{0} = \frac{2F - 4f}{k}, \quad E_{k} = \frac{F^{2} - 6fF + 8f^{2}}{k}$

（2）当 A 刚要离开墙时，设弹簧的伸长量为 $x$，以 A 为研究对象，由平衡条件得 $kx = f$

若 A 刚要离开墙壁时 B 的速度恰好等于零，这种情况下恒力为最小值 $F_{\min}$，从弹簧恢复原长到 A 刚要离开墙的过程中，以 B 和弹簧为研究对象，由能量守恒得

$E_{k} = \frac{1}{2}kx^{2} + fx$

结合第（1）问结果可知

$F_{\min} = \left(3 \pm \frac{\sqrt{10}}{2}\right)f$

根据题意舍去 $F_{\min} = \left(3 - \frac{\sqrt{10}}{2}\right)f$，所以恒力的最小值为

$F_{\min} = \left(3 + \frac{\sqrt{10}}{2}\right)f$

（3）从 B、C 分离到 B 停止运动，设 B 的路程为 $x_{B}$，C 的位移为 $x_{C}$，以 B 为研究对象，由动能定理得

$-W - fx_{B} = 0 - E_{k}$

以 C 为研究对象，由动能定理得

$-fx_{C} = 0 - E_{k}$

由 B、C 的运动关系得

$x_{B} > x_{C} - x_{BC}$

联立可知

$W < fx_{BC}$

（4）小物块 B、C 向左运动过程中，由动能定理得

$5fx_{1} - 2fx_{1} - \frac{1}{2}kx_{1}^{2} = 0$

解得撤去恒力瞬间弹簧弹力为

$kx_{1} = 6f$

则坐标原点的加速度为

$a_{1} = \frac{kx_{1} - 2f}{2m} = \frac{6f - 2f}{2m} = \frac{2f}{m}$

之后 C 开始向右运动过程（B、C 系统未脱离弹簧）加速度为

$a = \frac{kx - 2f}{2m}$

可知加速度随位移 $x$ 为线性关系，随着弹簧逐渐恢复原长，$x$ 减小，$a$ 减小，弹簧恢复原长时，B 和 C 分离，之后 C 只受地面的滑动摩擦力，加速度为

$a_{2} = -\frac{f}{m}$

负号表示 C 的加速度方向水平向左；从撤去恒力之后到弹簧恢复原长，以 B、C 为研究对象，由动能定理得

$\frac{1}{2}kx_{1}^{2} - 2fx_{1} = \frac{1}{2} \cdot 2mv^{2}$

脱离弹簧瞬间后 C 速度为 $v$，之后 C 受到滑动摩擦力减速至 0，由能量守恒得

$fx_{2} = \frac{1}{2}mv^{2}$

解得脱离弹簧后，C 运动的距离为

$x_{2} = \frac{1}{2}x_{1}$

则 C 最后停止的位移为

$x_{1} + x_{2} = \frac{3}{2}x_{1} = \frac{3}{2} \cdot \frac{6f}{k} = \frac{9f}{k}$

所以 C 向右运动的图象为

\onepicture[w=6.0cm]{figs/18_an.png}
}

\end{enumerate}

\end{document}
